\section{A minimum on open fibres and a compatible positivity source}
\label{in:minimality}

We work in the remaining case of Proposition~\ref{in:group-reduction}.
Thus $X$ is smooth connected projective, the support of $D$ is a
possibly empty disjoint union of smooth order-two divisors, and its
root orbifold is special.  There is a morphism
\begin{equation}\label{in:reduced-sequence}
 \alpha:X\longrightarrow A,
 \qquad
 1\longrightarrow K\longrightarrow G\longrightarrow
 \Lambda=\pi_1(A)\longrightarrow1,
\end{equation}
where $\alpha$ is surjective with connected fibres, the ordinary map
$\pi_1(X)\to\Lambda$ is an isomorphism, and $K$ is infinite, finitely
generated, and has no proper finite-index subgroup.  Every quotient of
$K$ has only trivial finite-dimensional complex linear images.  For a
general smooth fibre $P_a=\alpha^{-1}(a)$, the image of
$\pi_1(P_a\setminus D)$ in $G$ is exactly $K$.

\subsection{Descent of a reduction from an auxiliary manifold}

For clarity we first prove the descent statement used twice below.  If
$P$ is smooth projective, $P^\circ\subset P$ is a dense open, and
$\theta:\pi_1(P^\circ)\to E$ is a homomorphism, the image assigned
to a general fibre of a rational fibration of $P$ is computed on a
dense smooth open of a resolution which maps to $P^\circ$.  Removing
further proper algebraic subsets does not change that image, by
Lemma~\ref{in:topology}.  We call it the \emph{open fibre image}.

\begin{lemma}\label{in:descended-reduction}
Let $d:T\to P$ be a projective generically finite surjective map of
connected smooth projective varieties.  Suppose that the function-field
extension is Galois with finite group $J$, acting on $T$ over $P$.
Suppose $\theta:\pi_1(P^\circ)\to E$ induces a homomorphism
$\rho:\pi_1(T)\to E$, compatible on a dense open of $T$ mapping
to $P^\circ$.

There is a rational fibration $c:P\dashrightarrow Z$, with a smooth
projective base model, such that
\[
 \dim Z=\operatorname{gdim}(T,\rho)
\]
and the open group of its general fibre has finite image under
$\theta$.  The reduction of $(T,\rho)$ maps generically finitely to
$Z$, compatibly with $d$ and $c$.
\end{lemma}

\begin{proof}
Apply Assumption~\ref{in:input-reduction} to $(T,\rho)$ and write
$v:T\dashrightarrow Z_T$ for its reduction.  The base can be taken
projective.  Here is the usual justification in this setting.  Resolve
the graph of the meromorphic map from the projective source to a
compact K\"ahler model of the base.  Cutting the source by general
hyperplanes gives a projective subvariety of the base dimension mapping
generically finitely onto that base.  Under a proper generically finite
map the fields of meromorphic functions are a finite extension: after
Stein factorization, symmetric functions over the finite sheets give
the coefficients of the algebraic equations.  The base therefore has
full algebraic dimension.  It is Moishezon and admits a projective
birational model; resolving gives the asserted smooth projective model.

The finite-image test defining $v$ is invariant under $J$.  Indeed,
take a common $J$-stable dense open on which the compatibility with
$\theta$ holds.  The maps of this open to $P^\circ$ agree after deck
transformation.  With paths from a chosen point to its translates,
the induced homomorphisms differ by conjugation in $E$.  Images of
compact subvariety groups can be computed on dense opens of their
resolutions, because $\rho$ extends to $\pi_1(T)$.  Thus a deck
translate of a general finite-image fibre is again a finite-image
subvariety.  The contraction property makes it a fibre for the same
reduction, up to birational base equivalence.  Applying the argument
to inverse deck transformations gives invariance in both directions.

Consequently the relatively algebraically closed field
$F=\mathbb C(Z_T)\subset\mathbb C(T)$ is $J$-stable.  Its invariant
subfield
\[
 F^J\subset\mathbb C(T)^J=\mathbb C(P)
\]
is relatively algebraically closed in $\mathbb C(P)$.  To check this,
an element of $\mathbb C(P)$ algebraic over $F^J$ is algebraic over
$F$, hence belongs to $F$, and is $J$-invariant.  The field $F^J$
therefore defines a rational fibration $c:P\dashrightarrow Z$ with
connected general fibres.  The extension $F/F^J$ is finite, so the
base dimensions agree and $Z_T\dashrightarrow Z$ is generically
finite.

For general base points, a general fibre of $v$ maps dominantly and
generically finitely onto a general fibre of $c$: the source and base
maps in the commutative diagram are generically finite, and the two
fibre dimensions are equal.  On a smooth compact resolution the
upstairs fibre group has finite $\rho$-image.  The same is true on
any dense smooth open of that resolution, since its map to $T$ is
holomorphic and $\rho$ is defined on the compact group.  On suitable
opens the generically finite map has finite-index image on the
fundamental group of the downstairs fibre.  A group with a finite-index
subgroup of finite image itself has finite image.  Compatibility with
$\theta$ proves the assertion about the open fibre downstairs.
\end{proof}

For $P=X$, or $P=P_a$ with its smooth restricted boundary, the
auxiliary construction of Lemma~\ref{in:auxiliary} meets this lemma's
hypotheses for every quotient of the relevant open group image in
$G$.  Even pullback multiplicities, including on exceptional divisors,
are what extend the open group homomorphism to the compact manifold
$T$.

\subsection{Good fibres and a covering family}

Set $n=\dim X$ and $s=n-\dim A$.  We require a fixed notion of good
Albanese parameter which will work for all the fibrations subsequently
considered.

\begin{lemma}\label{in:cycle-families}
There is a countable collection of diagrams
\[
 h_i:V_i\longrightarrow T_i,
 \qquad j_i:V_i\longrightarrow X,
\]
such that $T_i$ is smooth connected quasi-projective, $h_i$ is smooth
proper with irreducible fibres, $j_i^{-1}D$ has relative simple normal
crossing support or is empty, and every compact irreducible subvariety
of $X$ not contained in $D$ is the image, birationally, of a fibre in
one of the families.

Outside a countable union of proper closed subsets of $A$, the fibre
$P_a$ is smooth, its boundary is smooth and disjoint, and no closure
$\overline{j_i(V_i)}\ne X$ contains $P_a$.
\end{lemma}

\begin{proof}
Use the countably many projective Hilbert schemes of $X$ and their
universal families.  The loci parametrizing reduced irreducible
subvarieties not contained in $D$ are constructible.  On a dense smooth
open of each irreducible parameter piece, resolve the reduced incidence
component dominating that piece, together with the inverse image of
$D$.  Characteristic-zero resolution and generic smoothness, applied
also to the finitely many boundary strata, give a smooth proper family
with relative normal crossings after further shrinking.  Its fibres
map birationally to the members: this holds on the generic fibre, and
the locus where the resolution is an isomorphism meets the members in
dense opens after shrinking.  The resolved fibres are geometrically
integral on that open.  Apply the same procedure to the complementary
lower-dimensional parameter pieces.  Noetherian induction, for each
Hilbert scheme, gives the required countable collection and loses no
member.  These are the usual projective parameter and resolution
constructions; see \cite{Hironaka,Stacks}.

Shrink $A$ so that $P_a$ and its boundary have the asserted smoothness
properties.  For any proper closed $Z_i=\overline{j_i(V_i)}$ dominating
$A$, the general fibre of $Z_i\to A$ has dimension at most $s-1$.
If $Z_i$ does not dominate $A$, its image is proper since $X\to A$
is proper.  In either case, outside a proper closed subset of $A$,
$Z_i$ cannot contain the irreducible smooth fibre $P_a$.  Take all
these countably many exclusions.  When $A$ is a point, a proper
$Z_i$ simply does not contain $X$.
\end{proof}

Call parameters satisfying Lemma~\ref{in:cycle-families} and the
open-fibre image conclusion $\operatorname{im}\pi_1(P_a\setminus D)=K$
\emph{good}.  They exist.  Moreover a good $P_a$ has points outside
all the nondominating evaluation closures: a countable union of
proper complex algebraic subsets does not cover a connected smooth
complex projective variety.

For every rational fibration $b:P_a\dashrightarrow V$ on a good
fibre, let $C_b\le K$ be its general open fibre image, computed in
$X\setminus D$.  Lemma~\ref{in:open-images} makes $C_b$ normal in
$K$.  Define
\begin{equation}\label{in:dimension-minimum}
 \ell=\min\{\dim V: a\text{ good},\quad
 b:P_a\dashrightarrow V\text{ a fibration},\quad K/C_b\text{ infinite}\}.
\end{equation}
The identity fibration is admissible, while a map to a point has
$C_b=K$.  Since $K$ is infinite and is an open fibre image, $s>0$,
and
\begin{equation}\label{in:ell-range}
 1\le\ell\le s.
\end{equation}
The minimum is over integers and is attained.  After any torus
isogeny pullback, a fibre over a lift of a good parameter is identified
with the old fibre, including its complement and its group images in
the same $K$.  Thus \eqref{in:dimension-minimum} remains a lower bound
for every such fibre and every further isogeny.  No countability of
the set of fibrations is needed for this definition.

Choose an attaining $b:P_a\dashrightarrow V$ and put $C=C_b$.
A general fibre image is a compact irreducible subvariety of dimension
$s-\ell$.  Choose it through a point outside all nondominating
evaluation closures, outside $D$, and in the good locus of the
resolved fibration.  It occurs in a family of
Lemma~\ref{in:cycle-families} whose evaluation is dominant.  Its open
group image in that family is $C$, since the two smooth birational
models have a common dense open.

The image of $\pi_1(V_i\setminus j_i^{-1}D)$ in $G$ normalizes $C$.
Indeed $h_i$, with its relative normal crossing boundary removed, is
a differentiably locally trivial family of connected opens; transport
identifies its fibre group images up to conjugation.  That total-space
image has finite index in $G$, by dominance of the evaluation and
Lemma~\ref{in:topology}.  The normalizer $G_1=N_G(C)$ therefore has
finite index and contains $K$.  It is the inverse image of a sublattice
of $\Lambda$.  Pass to the corresponding torus isogeny and its
connected coarse \'etale pullback of $X$.

The entire incidence family lifts to this pullback.  The open lifting
criterion holds by the preceding normalizer condition; composing to
$\pi_1(X)=\Lambda$ and using surjectivity from the incidence open to
$\pi_1(V_i)$ proves the lifting criterion for $V_i$ itself.  A
one-sheeted component of the finite \'etale pullback gives an algebraic
lift.  The lifted evaluation is dominant and its fibres still map
birationally to their images.  Transport now takes place in $G_1$,
where $C$ is normal, so all their open group images are $C$ with
compatible choices.  Relabel the resulting data as $X,A,G$.
We have obtained
\begin{equation}\label{in:normal-C}
 C\triangleleft G,\qquad C\subset K,\qquad K/C\text{ infinite},
\end{equation}
and a covering family of compact subvarieties of dimension $s-\ell$
whose open group images lie in $C$.

\subsection{The intermediate fibration and its quotients}

\begin{proposition}\label{in:intermediate}
After the preceding isogeny there is a factorization by fibrations
\[
 X\overset{c}{\dashrightarrow}Z\overset{g}{\longrightarrow}A,
 \qquad \alpha=g\circ c,
\]
with $Z$ smooth projective, such that the general open $c$-fibre image
$B_0$ is normal in $G$, is contained in $K$, and satisfies
\[
 \dim Z-\dim A=\ell>0,
 \qquad Q=K/B_0\text{ infinite}.
\]
Writing $R=G/B_0$, one has an exact sequence
\begin{equation}\label{in:RQ-sequence}
 1\longrightarrow Q\longrightarrow R\longrightarrow\Lambda
 \longrightarrow1.
\end{equation}
Every quotient of $Q$ has no proper finite-index subgroup and has only
trivial finite-dimensional complex linear images.
\end{proposition}

\begin{proof}
Choose an auxiliary $d:T\to X$ from Lemma~\ref{in:auxiliary}, and
apply Lemma~\ref{in:descended-reduction} to the quotient $G/C$.
Let $v:T\dashrightarrow Z_T$ be the manifold reduction and
$c:X\dashrightarrow Z$ its descended fibration.

The lifted covering family forces
$\dim Z_T\le n-(s-\ell)=\dim A+\ell$.  To verify this directly,
work over a dense open where $d$ is finite \'etale and the evaluation
family covers.  A member through the projection of a point of $T$
has a component of its inverse image through that point, of the same
dimension $s-\ell$, dominating the member.  On a resolution of this
compact component, its group image under $\pi_1(T)\to G/C$ is
trivial: compute on a common dense open mapping to the resolved member,
whose open image lies in $C$.  Surjectivity from that open to the
compact resolution group is applicable because the homomorphism
extends to $\pi_1(T)$.  These compact subvarieties pass through points
in a dense open.  At very general such points the contraction property
of Assumption~\ref{in:input-reduction} applies.  Its general fibres
therefore have dimension at least $s-\ell$.

The map $T\to A$ is constant on a general resolved fibre of $v$.
That fibre has finite group image in $G/C$ and thus finite image in
the torsion-free lattice $\Lambda$.  Its map to $A$ lifts to the
vector-space universal cover and is constant by compactness.  It
follows that $T\to A$ factors rationally through $Z_T$.  The factor
is invariant under the deck group, so it descends to $Z$.  Resolve
to make $g:Z\to A$ holomorphic.  Its general fibres are connected:
$\mathbb C(A)$ is relatively algebraically closed in $\mathbb C(X)$
by connectedness of the general $\alpha$-fibre, and hence also in
$\mathbb C(Z)$.

Let $B_0$ be the general open $c$-fibre image.  It is normal in $G$
by Lemma~\ref{in:open-images}, and lies in $K$ because $c$ is over
$A$.  Lemma~\ref{in:descended-reduction} says that its image in
$G/C$ is finite.  If $K/B_0$ were finite, the image of $K$ in $G/C$
would consequently be finite, contradicting \eqref{in:normal-C}.
Thus $Q$ is infinite.

For a very general good $a$, restrict the diagram to $P_a$.  The
result is a rational fibration onto the corresponding smooth connected
$g$-fibre, with general open fibre image $B_0$ in $K$.  This follows
on the common resolved fibration locus: the fibres are the same
general $c$-fibres, and deleting additional algebraic subsets preserves
their images.  Normality of $B_0$ in $G$ removes the ambiguity from
transporting paths.  The quotient $K/B_0$ is infinite, so the minimum
\eqref{in:dimension-minimum} gives $\dim Z-\dim A\ge\ell$.
The opposite inequality was already proved.  The exact sequence and
the remaining group properties follow from $B_0\subset K$ and from
the corresponding properties of $K$.
\end{proof}

We record carefully the effect of further isogenies.  If $A'\to A$
is one, the pullbacks of $X$ and a smooth model of $Z$ are connected.
For $X$ this follows from $\pi_1(X)\simeq\Lambda$; for $Z$ it
follows because a proper fibration with connected general fibres
induces a surjection on $\pi_1(A)$, as seen first over the smooth
fibration open.  The pulled-back $c$-fibres identify with the old
ones over chosen lifts of base points.  Thus their open group image,
under the inclusion $G'\le G$, is the same $B_0$, and
\begin{equation}\label{in:isogeny-groups}
 R'=G'/B_0\hookrightarrow R,
 \qquad Q=K/B_0\text{ is unchanged}.
\end{equation}
The relative dimension, the good-fibre minimum, and the source
specialness are unchanged as well.  We shall use these facts for
arbitrary further torus isogenies.

\subsection{Actual inertia and fibrewise bigness}

The open-fibration construction gives a compatible surjection
\[
 \psi:\pi_1(Z^\circ)\twoheadrightarrow R.
\]
On a smooth projective model over $A$, set $m_L$ equal to the order
of the image of a divisorial meridian about a prime $L$, and set it
equal to one for primes meeting $Z^\circ$.  Lemma~\ref{in:actual-inertia}
shows that these orders are finite.  Put
\begin{equation}\label{in:true-boundary}
 \Delta_Z=\sum_L(1-m_L^{-1})L.
\end{equation}
Monodromy is transported between models on common opens.  Strict
transforms retain their orders.  Under an isogeny pullback the
meridians lift with winding number one and
\eqref{in:isogeny-groups} preserves their orders.

Call a model \emph{prepared} if it is smooth projective over $A$ and
the divisorial part of a deleted locus defining this monodromy has
simple normal crossings. A prepared model can be chosen above every
smooth projective model in use: take common resolutions, resolve the
complement of a common monodromy open on the base, and resolve the
source diagram and its deleted sets.  Generic smoothness for all
boundary strata provides the required smooth family of complements
over a smaller base open.  The construction works after every torus
isogeny.  The coefficients on newly appearing primes are always read
from their actual meridians.

\begin{proposition}\label{in:fibre-big}
On every prepared model just described, including those after any
further torus isogeny, a very general smooth fibre $Y$ of $g:Z\to A$
satisfies
\[
 \dim Y=\ell,
 \qquad K_Y+\Delta_Z|_Y\text{ is big}.
\]
\end{proposition}

\begin{proof}
Choose a very general parameter $a$ lying above a good parameter of
the original torus, and put $Y=g^{-1}(a)$. Smoothness, connectedness, and dimension
$\ell$ follow from Proposition~\ref{in:intermediate}.  Generic
smoothness for the finitely many deleted strata ensures that
$\Delta_Y=\Delta_Z|_Y$ has simple normal crossing support and finite
integral root coefficients.  Higher-codimension deletions remain of
higher codimension on a very general fibre or miss it.

The restriction of $\psi$ to the open part $Y^\circ$ has image in
$Q$, since its composite with the lattice quotient is induced by the
constant map $Y\to A$.  It is onto $Q$: a dense open of $P_a$ maps
compatibly to $Y^\circ$, and its group image in $G$ is $K$.
The meridians of deleted divisors on $Y$ are transverse restrictions
of the corresponding meridians on $Z$, with the same images up to
transport.  Their orders in $Q$ equal their orders in $R$ because
$Q\le R$.  Normal meridian generation therefore gives a surjection
\begin{equation}\label{in:q-positive}
 q:\pi_1^{\mathrm{orb}}(Y,\Delta_Y)\twoheadrightarrow Q.
\end{equation}

Apply Lemma~\ref{in:auxiliary} to the smooth projective fibre $P_a$
with its smooth disjoint or empty boundary.  It gives a connected
smooth projective Galois auxiliary map $d:S\to P_a$ and a compatible
homomorphism $\pi_1(S)\to K$.  This homomorphism has finite-index
image: on dense opens the generically finite map supplies a
finite-index subgroup of the open group of $P_a$, whose image is
$K$.  Since $K$ has no proper finite-index subgroup, its image is
all of $K$.  Composition with $K\twoheadrightarrow Q$ gives
$\nu:\pi_1(S)\twoheadrightarrow Q$.  The dominant meromorphic map
$S\dashrightarrow Y$ and $\nu$ are compatible with
\eqref{in:q-positive} on a dense open by the construction of
monodromy.

It remains to verify the last, quantified hypothesis of
Assumption~\ref{in:input-positive}.  There is only one finite-index
subgroup of $Q$, namely $Q$ itself, so use the identity cover of this
fixed $S$.  Let $\psi_0:Q\twoheadrightarrow H$ be any infinite
quotient.  Apply Lemma~\ref{in:descended-reduction} on $P_a$ to
$S\to P_a$ and the compatible composition to $H$.  It gives a
rational fibration of $P_a$ whose base dimension is
$\operatorname{gdim}(S,\psi_0\nu)$ and whose general open fibre
image $C'\triangleleft K$ has finite image in $H$.  The quotient
$K/C'$ is infinite, since $K$ surjects onto the infinite group $H$
and the image of $C'$ there is finite.  This is one of the fibrations
tested by \eqref{in:dimension-minimum}.  Hence
\[
 \operatorname{gdim}(S,\psi_0\nu)\ge\ell.
\]
This argument holds for every infinite quotient of $Q$ on the
\emph{same} $S$.

We have now checked all the hypotheses of
Assumption~\ref{in:input-positive}: $Y$ and $S$ are connected smooth
projective manifolds, the source dominates compatibly and surjects on
the required group, the boundary is finite integral SNC, $Q$ is
infinite with finite linear images, and the full quotient condition
holds.  It proves the desired bigness.  The proof applies separately
to every prepared model and every further isogeny, using
\eqref{in:isogeny-groups} and the unchanged lower bound
\eqref{in:dimension-minimum}.  No single parameter is required to be
very general simultaneously for all possible models.
\end{proof}
